
est constante sur $X$. Comme $f$ est projectif, il existe une immersion fermée $X\to P$, où $P$ est un fibré projectif sur $Y$, et comme $Y$ a un faisceau ample, on peut supposer que $P=\mathbf{P}^r_Y$. On a donc une factorisation :
\[
\begin{tikzcd}[column sep=large,row sep=large]
X \arrow[rr,"i"] \arrow[dr,"f"'] && \mathbf{P}^r_Y \arrow[dl,"g"]\\
& Y. &
\end{tikzcd}
\]
Comme $f$ est d'intersection complète, l'immersion $i$ est régulière (1.2).

D'après 1.7, $f$ est un morphisme parfait. De plus, $Y$ a un faisceau ample et par suite $X$ et $P$ également. On peut alors définir un homomorphisme additif :
\[
f_*:K^{\bullet}(X)\longrightarrow K^{\bullet}(Y)
\]
(IV 2.12), les $K^{\bullet}$ définis en termes de faisceaux localement libres de type fini, ou de complexes parfaits, étant identiques. On peut également définir des homomorphismes additifs :
\[
i_*:K^{\bullet}(X)\longrightarrow K^{\bullet}(P);
\qquad
 g_*:K^{\bullet}(P)\longrightarrow K^{\bullet}(Y),
\]
et on a : $f_*=g_*\circ i_*$.

Soit $d$ la dimension relative virtuelle de $X$ sur $Y$. La codimension de $i$ est alors $r-d$. D'après VII 4.6, l'homomorphisme $i_*$ donne après tensorisation par $\mathbb Q$ des inclusions :
\[
i_*\bigl(\Fil^kK^{\bullet}(X)_{\mathbb Q}\bigr)
\subset
\Fil^{k+r-d}K^{\bullet}(P)_{\mathbb Q}.
\]
D'après VI 5.8, l'homomorphisme $g_*$ donne des inclusions :
\[
g_*\bigl(\Fil^{k'}K^{\bullet}(P)\bigr)
\subset
\Fil^{k'-r}K^{\bullet}(Y).
\]
En composant les deux, on obtient donc:

